\section{Nonempty coefficient families and integral arithmetic realizations}
\label{sec:nonempty-realization}

\status{NR1--NR18 construct the actual all-prime coefficient family
\texttt{TwR/RingSpec}, its normalized characters, finite universal
anticyclotomic quotients, and readings at every prescribed coefficient
place. NR19--NR24 prove finite-field descent and integralization for
compact representations and the source's bounded family of row
representations, conditional on its stated automorphic Galois inputs.
The original number fields, rows, complex embeddings and host are
fixed throughout. This supplies nonempty twist indices and the exact
coefficient part of \texttt{SideB.cover}; it does not supply the
remaining geometric and automorphic form realizations.}

Fix the source datum $D$, its host CM field $E$, an algebraic closure
$\Omega/E$, its prescribed places on $\Omega$, and its character frame
$C$. Put
\[
 K=D.\mathrm{LL0}\subset L=D.\mathrm{what}_2.\mathrm{LL},\qquad
 \epsilon:L\hookrightarrow\mathbb C .                  \tag{NR.1}
\]
Both $K$ and $L$ are number fields. Whenever conjugate coefficient
embeddings are used, retain the literal hypothesis $\bar K=K$, with
conjugation induced from the fixed $L$ and $\epsilon$.
For a rational prime $\ell$, write $\Omega_\ell=\overline{\mathbb Q}_\ell$
and $C_\ell=\widehat{\Omega_\ell}$, with $|\ell|=\ell^{-1}$.
The algebraic closure $\Omega_\ell$ is not complete. A finite extension
always means a finite subextension of $\Omega_\ell/\mathbb Q_\ell$.
Its integer ring is the closed unit ball. We use the standard local
field facts that finite extensions are complete and their extended
absolute value is unique. The integer-ring properties needed here
are verified in NR3. The algebraic closedness of $C_\ell$, elementary
finite-dimensional nonarchimedean norm equivalence, the Baire theorem,
and finiteness of number-field ideal class groups are stated classical
inputs. Global Artin reciprocity, where used, has exactly the source's
normalization; it is not inferred from these local constructions.

\paragraph{NR1: embeddings realizing each actual finite prime.}
Let $q$ be a nonzero prime of $\mathcal O_K$, and let $\ell$ be the
characteristic of its finite residue field. Thus
$q\cap\mathbb Z=(\ell)$, with $\ell$ a rational prime. The completion
$K_q$ is a finite extension of $\mathbb Q_\ell$. Choose an embedding
$K_q\hookrightarrow\Omega_\ell$ over $\mathbb Q_\ell$ and restrict it
to $K$. Uniqueness of the extended absolute value gives an embedding
$J_q:K\hookrightarrow\Omega_\ell$ satisfying
\[
 a\in q\quad\Longleftrightarrow\quad |J_q(a)|<1
       \qquad(a\in\mathcal O_K).                      \tag{NR.2}
\]
This is the literal \texttt{exists\_embAt} condition, so the prescribed
family $J$ in \texttt{exists\_ringSpec} can be supplied, rather than
left as a new existence premise.

More generally, let $A$ be a valuation ring on $\Omega$ above a prime
of $E$ of residue characteristic $\ell$. There is an embedding
$j:\Omega\hookrightarrow\Omega_\ell$ such that
\[
                    x\in A\quad\Longleftrightarrow\quad |j(x)|\le1.
                                                               \tag{NR.3}
\]
One proof first chooses any embedding of the algebraic closure of
$\mathbb Q$ into $\Omega_\ell$. Its pulled-back valuation and $A$
restrict to the same $\ell$-adic valuation ring on $\mathbb Q$.
Extensions of a valuation to an algebraic normal extension are
conjugate: in each finite normal subextension this is transitivity on
primes above the given prime, and compactness of the inverse system
of the nonempty finite sets of conjugating automorphisms gives one
automorphism on the algebraic closure. Precomposing the first
embedding by its inverse proves (NR.3). This argument also explains
why a choice of prime must accompany a global algebraic embedding.

\paragraph{NR2: density and uniqueness of the normalized local norm.}
For every embedding $j:\Omega\hookrightarrow\Omega_\ell$, its image
is dense. Indeed, let $a\in\Omega_\ell$, choose a monic polynomial
$f\in\mathbb Q_\ell[X]$ with $f(a)=0$, and approximate its finitely
many coefficients by rational numbers. For any $\varepsilon>0$ this
gives a monic $g\in\mathbb Q[X]$ of the same degree $d$ with
$|g(a)|<\varepsilon^d$. All roots $b_1,\ldots,b_d$ of $g$ belong to
$j(\Omega)$, since this is an algebraically closed field containing
$\mathbb Q$. The equality
$|g(a)|=\prod_i|a-b_i|$ implies $|a-b_i|<\varepsilon$ for some $i$.

If $j,j':\Omega\hookrightarrow\Omega_\ell$ pull back the same
valuation ring, then $|j(x)|=|j'(x)|$ for every $x$. For $x\ne0$, test
membership of $\ell^{-m}x^n$ for every integer $m$ and positive integer
$n$. This compares $|j(x)|$ and $|j'(x)|$ with all numbers
$\ell^{m/n}$; these comparisons determine a positive real number.
Thus the map $j(x)\mapsto j'(x)$ is an isometry. Extend it and its
inverse to the complete field $C_\ell$. The extensions are inverse
field isometries, fix $\mathbb Q_\ell$ by density of $\mathbb Q$, and
preserve the elements algebraic over $\mathbb Q_\ell$. Restriction
therefore gives a unique
\[
             \tau\in\operatorname{Aut}_{\mathbb Q_\ell}(\Omega_\ell),
             \qquad \tau j=j'.                         \tag{NR.4}
\]
The completion is used in this proof; an extension theorem into the
incomplete field $\Omega_\ell$ has not been assumed.

\paragraph{NR3: the actual local decomposition bridge.}
For $j$ as in (NR.3), put $K_v=\mathbb Q_\ell(j(E))$.
A rational basis of $E$ gives finitely many algebraic generators,
so $K_v/\mathbb Q_\ell$ is finite. Its closure is the closure of
$j(E)$: the latter contains the closure $\mathbb Q_\ell$ of
$\mathbb Q$, and is a field containing the finitely many generators;
conversely the finite field $K_v$ is closed.
An element $\sigma$ of the decomposition group $D_A\subset G_E$
preserves $A$. NR2 extends $j\sigma j^{-1}$ uniquely to an
automorphism of $\Omega_\ell$ fixing $K_v$. Conversely an automorphism
fixing $K_v$ preserves the set $j(\Omega)$ of elements algebraic over
$\mathbb Q$ and the $\ell$-adic norm. Its restriction gives an element
of $D_A$. These constructions are inverse and yield
\[
       D_A\simeq\operatorname{Gal}(\Omega_\ell/K_v),
       \qquad j(\sigma x)=\tau_\sigma(jx).              \tag{NR.5}
\]
The inverse map is continuous: the condition of fixing finitely many
elements of $\Omega$ pulls back to the open stabilizer of finitely
many algebraic elements of $\Omega_\ell$. Both groups are compact
Hausdorff profinite groups, so a continuous bijection is a
homeomorphism. This proves the required topology in both directions.

Choose a finite normal extension $L_v/\mathbb Q_\ell$ containing
$K_v$. Every $\mathbb Q_\ell$-embedding of $K_v$ lands in $L_v$.
For every finite extension $K_0\supset L_v$, the ring
$\mathcal O_{K_0}$, the map $j$, and (NR.5) therefore give the
literal \texttt{LocBridge}, including all its coefficient embeddings.
The equivalence (NR.3) is its \texttt{br\_ind} clause. This is a
bridge for the prescribed place $A$, not a replacement place.

Here are the needed integer-ring facts, to make the base-ring
construction explicit as well. If $K_0/\mathbb Q_\ell$ has degree
$d$, the field norm gives $|x|^d\in\ell^{\mathbb Z}$ for $x\ne0$.
Choose an element of least positive value. Every nonzero element
of the unit ball is a unit times a nonnegative power of it, by
repeated division, so the unit ball $O_0$ is a DVR. Relative to
a $\mathbb Q_\ell$-basis of $K_0$, norm equivalence bounds all
coordinates of its unit ball. Thus $O_0$ is a $\mathbb Z_\ell$-submodule
of a finite fractional $\mathbb Z_\ell$-lattice; it is finite and
torsion-free, hence free. It is consequently integral over
$\mathbb Z_\ell$. Conversely a monic integral equation forbids
norm greater than one, so this unit ball is exactly the integral
closure used in the source. Multiplication by a sufficiently
large power of $\ell$ puts any element of $K_0$ into $O_0$, proving
$\operatorname{Frac}(O_0)=K_0$. Its residue field is a quotient of
the finite ring $O_0/\ell O_0$, and is finite. Powers of a
uniformizer are precisely successively smaller closed norm
balls, so the topology is adic, and closedness in the complete
field gives adic completeness. In the bridge the coefficient
isomorphism is the identity of $O_0$, hence is continuous in
both directions and respects $\mathbb Z_\ell$.

\paragraph{NR4: the value group of the unramified algebraic field.}
Let $U_\ell\subset\Omega_\ell$ be the field generated by all roots of
unity of order prime to $\ell$, the source's maximal unramified
extension. For $(m,\ell)=1$, set
$R_m=\mathbb Z_\ell[\mu_m]$. Products of $m$th roots are again $m$th
roots, so $R_m$ is their finite $\mathbb Z_\ell$-span. It is a compact,
and hence closed, subset of $\Omega_\ell$, and all its elements have
norm at most one. Put $N=\ell^{\varphi(m)}>1$. Euler's congruence
gives $\zeta^N=\zeta$ for $\zeta^m=1$, and reduction modulo $\ell$
gives $a^N\equiv a$ for $a\in\mathbb Z_\ell$. Applying Frobenius in
the quotient $R_m/\ell R_m$ therefore gives
\[
                   y^N-y\in\ell R_m\qquad(y\in R_m).  \tag{NR.6}
\]
If $|y|<1$, iteration shows
$\ell^{-1}(y-y^{N^s})\in R_m$ for every $s$. These elements tend
to $y/\ell$; closedness gives $y/\ell\in R_m$.
For $y\ne0$, repeat this division while the norm is less than one.
The process must terminate, since the norm is multiplied by $\ell$
at each step and elements of $R_m$ have norm at most one. Consequently
$|y|=\ell^{-a}$ for an integer $a\ge0$.

Every element of $U_\ell$ is a quotient of two elements of some
$R_m$. To check this without a finiteness assumption on $U_\ell$,
take the finitely many roots occurring in a rational expression,
replace their orders by their least common multiple, and clear the
finitely many $\mathbb Q_\ell$ denominators. Hence
\[
                        |U_\ell^\times|=\ell^{\mathbb Z}.
                                                               \tag{NR.7}
\]
This proves the discrete norm assertion used below directly from
the prime-to-$\ell$ roots.

\paragraph{NR5: completion preserves a discrete valuation ring.}
Let $L_0/\mathbb Q_\ell$ be finite and put $F=L_0U_\ell$.
The extension $F/U_\ell$ has finite degree $d>0$. Every embedding
of $F$ over $U_\ell$ preserves the $\ell$-adic norm, because it fixes
$\mathbb Q_\ell$ and the extended norm on an algebraic extension is
unique. The field norm thus satisfies
\[
 |N_{F/U_\ell}(x)|=|x|^d,\qquad
 |x|^d\in\ell^{\mathbb Z}\quad(x\in F^\times).         \tag{NR.8}
\]
The positive additive values of $F$ lie in $d^{-1}\mathbb Z$ and
form a nonempty set, since it contains the value of $\ell$.
Choose $\pi\in F$ having its least positive value. Equivalently,
$0<|\pi|<1$ is the largest norm of an element of norm below one.

Let $\widehat F$ be the closure of $F$ in $C_\ell$ and put
$O=\{x\in\widehat F:|x|\le1\}$. The norm gap survives completion:
the closed set
$\{x:|x|\le|\pi|\}\cup\{x:|x|\ge1\}$ contains $F$, hence
contains its closure. Dividing repeatedly by $\pi$ now shows that
every nonzero $x\in O$ has a unique form $u\pi^a$ with $u\in O^\times$
and $a\ge0$. For an ideal, the least value of a nonzero element
generates it. Thus $O$ is a DVR with maximal ideal $(\pi)$, and
\[
 x\in(\pi)^a\quad\Longleftrightarrow\quad |x|\le|\pi|^a.
                                                               \tag{NR.9}
\]
Its induced topology is exactly its maximal-ideal adic topology.
The field $\widehat F$ and its closed unit ball are complete, so $O$
is also adically complete.

Its residue field is algebraically closed. The residue field of
the algebraically closed valued field $C_\ell$ is algebraically
closed: lift a monic residue polynomial, take a root in $C_\ell$,
and use its monic equation to see that every root is integral.
The map $O/(\pi)\to\overline{\mathbb F}_\ell$ is injective.
It is onto because every nonzero residue is represented by a
prime-to-$\ell$ root of unity in $U_\ell$. Explicitly, approximate
a norm-one element of $C_\ell$ to distance less than one by an
algebraic element; in a finite residue field its nonzero reduction
has order prime to $\ell$, and the simple root of
$X^{\ell^f-1}-1$ lifts to a root of unity with that reduction
(equivalently, use the residue-root lifting of LS8).
Finally $O$ is the closure of the unit ball of $F$: sufficiently
close approximants to an element of norm at most one also have
norm at most one. These are all the assertions of
\texttt{exists\_integers\_completion}; $O$ is not the full,
nondiscrete valuation ring of $C_\ell$.

\paragraph{NR6: bounded coefficients after one algebraic adjunction.}
Let $F\subset\Omega_\ell$ contain $\mathbb Q_\ell$, let
$\widehat F\subset C_\ell$ be its closure, and let $z\in\Omega_\ell$
be algebraic over $\mathbb Q_\ell$. If a monic polynomial of degree
$n>0$ over $\mathbb Q_\ell$ kills $z$, the image of
\[
 T:\widehat F^n\longrightarrow C_\ell,
       \qquad (a_0,\ldots,a_{n-1})\longmapsto
                         \sum_{a<n}a_a z^a             \tag{NR.10}
\]
is $\widehat F[z]=\widehat F(z)$. Indeed, polynomial division spans
the algebra by these powers, and a finite-dimensional domain over a
field is a field. This image is a finite-dimensional normed space
over the complete field $\widehat F$, so it is complete and closed.
It contains $F(z)$, and every displayed linear combination is a
limit of combinations with coefficients in $F$. It follows that
\[
                     \overline{F(z)}=\widehat F(z).
                                                               \tag{NR.11}
\]
Choose a vector-space complement to $\ker T$. On that complement,
$T$ is an isomorphism of finite-dimensional normed spaces, whose
inverse is bounded. There is consequently a constant $C>0$ such
that every $x\in\overline{F(z)}$ with $|x|\le1$ admits coefficients
in (NR.10) satisfying $|a_a|\le C$. Powers of $z$ need not remain
linearly independent over $\widehat F$; the complement is what
justifies the bound when the degree drops after completion.

\paragraph{NR7: the finite cyclotomic coefficient tower.}
For finite $K_0,L_0\subset\Omega_\ell$, let
\[
 F=K_0L_0U_\ell,\quad O=\mathcal O_{\widehat F},\quad
 F_k=F(\zeta_{\ell^k}),\quad O_k=\mathcal O_{\widehat{F_k}},
                                                               \tag{NR.12}
\]
where $\zeta_{\ell^k}$ is primitive and $\zeta_1=1$.
NR5 applies to both $F$ and $F_k$, so $O,O_k$ are complete DVRs
with algebraically closed residue fields and their exact adic
topologies. Their inclusions are injective and $\ell$ is in their
maximal ideals.

Apply NR6 with $z=\zeta_{\ell^k}$. Choose $N$ so large that
$|\ell|^N C\le1$. For every $x\in O_k$,
\[
              \ell^N x\in
                O+O\zeta_{\ell^k}+\cdots+O\zeta_{\ell^k}^{n-1}.
                                                               \tag{NR.13}
\]
Thus $O_k$ is an $O$-submodule of a finite fractional $O$-module.
Since the DVR $O$ is Noetherian, $O_k$ is a finite $O$-module.
For any $x\in\operatorname{Frac}(O_k)$, multiply first by a
sufficient power of $\ell$ to make it integral and then use
(NR.13). This proves the exact generation statement
\[
 \exists d\in O\setminus\{0\}\ \exists p\in O[X],
                         \qquad dx=p(\zeta_{\ell^k}).  \tag{NR.14}
\]
The root has norm one and is a unit of exact order $\ell^k$.
For $k=0$ the two rings agree. These statements supply
\texttt{Ok\_fin}, \texttt{zk\_prim}, and \texttt{Ok\_gen}; the
stronger, unnecessary assertion $O_k=O[\zeta_{\ell^k}]$ has not
been used.

\paragraph{NR8: principal units cannot carry prime-to-residue torsion.}
Let $v$ be a finite place of $E$ of residue characteristic $p\ne\ell$,
and let $\chi:E_v^\times\to B$ be a continuous character with open
kernel, where every element of the finite group $B$ has
$\ell$-power order. If $u\in1+\mathfrak m_v$, then
$u^{p^a}\to1$. For completeness, if $u=1+x$ with
$\operatorname{ord}_v(x)\ge1$, the binomial theorem gives
\[
 \operatorname{ord}_v((1+x)^p-1)
          \ge \operatorname{ord}_v(x)+1.              \tag{NR.15}
\]
The middle binomial coefficients are divisible by $p$, and the
last term has value $p\operatorname{ord}_v(x)$, so this also covers
$p=2$. Iteration proves convergence. The open kernel eventually
contains $u^{p^a}$; hence $\chi(u)^{p^a}=1$. Its order is both an
$\ell$-power and a divisor of $p^a$, so $\chi(u)=1$.

On an infinite complex component every homomorphism to a finite
group is trivial: if the target has exponent $M$, write any
nonzero complex number as an $M$th power. Thus an idele character
as above, unramified outside a finite set $Q$ and also at every
prime over $\ell$, kills
\[
 U_Q=E_\infty^\times\times
       \prod_{v\in Q}(1+\mathfrak m_v)
                      \times\prod_{v\notin Q}\mathcal O_{E_v}^\times.
                                                               \tag{NR.16}
\]
At $v\in Q$ over $\ell$, its assumed unramifiedness kills all units;
at the other primes of $Q$, (NR.15) applies. To pass from local
components to the full product, approximate a unit idele by the
net of its finite-support components. The character kernel is
closed because it is an open subgroup, so it contains the limit.

\paragraph{NR9: a finite common idelic quotient.}
Let $U=E_\infty^\times\prod_v\mathcal O_{E_v}^\times$.
The quotient $U/U_Q$ injects into
$\prod_{v\in Q}\kappa(v)^\times$, and is finite. The valuation
map from ideles to fractional ideals is surjective, has kernel
$U$, and carries the diagonal $E^\times$ to principal ideals.
Consequently
\[
 \mathbb I_E/(E^\times U)\simeq\operatorname{Cl}(E),\qquad
                  [\mathbb I_E:E^\times U_Q]<\infty.  \tag{NR.17}
\]
Here finiteness of $\operatorname{Cl}(E)$ is the classical ideal
class theorem. No uniform bound for its order is required.
Equivalently, any subgroup of $\mathbb I_E$ containing $E^\times$
and $U_Q$ has finite index: first use the finite residue quotient
inside $U$, then the ideal class quotient. This is exactly the
common finiteness statement needed for all the characters in
NR10; their conductor exponents do not have to be bounded separately.

\paragraph{NR10: the universal kernel is finite index and open.}
Fix the prescribed conjugation $c$ of $\Omega$, a rational prime
$\ell$, and a finite set $Q$ of primes of $E$. Let $\mathcal K_Q$
be the kernels of all homomorphisms from $G_E$ to finite abelian
$\ell$-primary groups which have open kernel, satisfy
$\phi(cgc^{-1})=\phi(g)^{-1}$, and are unramified at every prime
outside $Q$ \emph{and at every prime over $\ell$}. Define
\[
                         H_Q=\bigcap_{H\in\mathcal K_Q}H.
                                                               \tag{NR.18}
\]
The trivial character ensures that this is an intersection of a
nonempty family. Each kernel is closed and contains the closure
of the commutator subgroup. Thus $G_E/H_Q$ is an abelian quotient
of $G_E^{\mathrm{ab}}$.

Use the supplied Global Artin reciprocity map
$a_E:\mathbb I_E\to G_E^{\mathrm{ab}}$. Its exact source quotient
map is surjective and kills diagonal elements. Its local unit
law identifies the image of local units with inertia. By NR8,
the composite of $a_E$ with every character defining
$\mathcal K_Q$ kills $U_Q$. Hence the composite
$\mathbb I_E\to G_E/H_Q$ kills $E^\times U_Q$ and is surjective.
NR9 makes $G_E/H_Q$ finite.

There are now only finitely many distinct subgroups $H$ in
$\mathcal K_Q$, because each contains $H_Q$ and corresponds to a
subgroup of the finite quotient. A finite intersection of such
kernels is again of the required kind: use the product of the
finitely many characters, whose target remains finite abelian
and $\ell$-primary. Therefore $H_Q$ itself is one of the kernels.
In particular it is open, and its quotient is $\ell$-primary.
This last argument is essential: an arbitrary infinite
intersection of open subgroups need not be open.

\paragraph{NR11: universal anticyclotomic quotients and restriction.}
Set
\[
     \operatorname{RCG}_{\ell,Q}=G_E/H_Q,
             \qquad a_{\ell,Q}:G_E\twoheadrightarrow
                                      \operatorname{RCG}_{\ell,Q}.
                                                               \tag{NR.19}
\]
NR10 proves finiteness, an open kernel and $\ell$-power order for
every element. The definitions prove the anticyclotomic law and
unramifiedness outside $Q$, as well as unramifiedness at primes
over $\ell$. If $Q$ contains no prime over $\ell$, every finite
abelian $\ell$-primary anticyclotomic character with open kernel
and ramification contained in $Q$ occurs in (NR.18). It therefore
factors uniquely through $a_{\ell,Q}$. If $Q'\subset Q$, the family
of permitted characters for $Q'$ is smaller, so
$H_Q\subset H_{Q'}$. The quotient map gives
\[
 \operatorname{res}_{Q,Q'}:\operatorname{RCG}_{\ell,Q}
      \longrightarrow\operatorname{RCG}_{\ell,Q'},\qquad
 \operatorname{res}_{Q,Q'}a_{\ell,Q}=a_{\ell,Q'}.        \tag{NR.20}
\]
These maps are compatible under further restriction by surjectivity
of $a_{\ell,Q}$.
This proves the full \texttt{exists\_ringClass\_host} supplier.
The restriction that $Q$ avoid $\ell$ in \texttt{art\_univ} is
retained exactly. The quotient exists also for other $Q$, but it
still kills inertia at $\ell$; it is not asserted to classify
arbitrarily ramified $\ell$-adic characters. Nor does universality
assert that this finite quotient is nontrivial or that a character
ramifies at every prime in $Q$.

\paragraph{NR12: the lower bound is chosen before the coefficient rings.}
For every $(\ell,j)$, where $j:K\hookrightarrow\Omega_\ell$, define
$c_0(\ell,j)$ to be the compositum of
\[
 \mathbb Q_\ell(j(K)),\qquad
 \mathbb Q_\ell(\mu_\ell),\qquad
                 L_v\quad(v\mid\ell\text{ in }E),    \tag{NR.21}
\]
where the $L_v$ are the finite normal fields from NR3 for the
prescribed places on $\Omega$. There are finitely many primes
above $\ell$, so this is a finite extension. It is chosen before
any of the fields denoted $b_0,b_1$ in the source. For every
$K_0\supset c_0(\ell,j)$ it contains $j(K)$, the $\ell$th roots,
and the coefficients of every required local bridge. This is
the literal \texttt{exists\_coefficientField\_lower\_bounds},
with its order of quantifiers unchanged.

We also need compatible complex readings. Each of $\Omega_\ell$
and $C_\ell$ is an algebraically closed field of characteristic
zero and cardinality continuum. The lower cardinal bound comes
from $\mathbb Q_\ell$; the upper bound for the algebraic closure
comes from counting polynomials, and for its metric completion
from counting sequences. An algebraically closed field of this
uncountable cardinality has transcendence degree continuum over
the countable field $K$. Extend the prescribed isomorphism
$j(K)\to\epsilon(K)$ to their algebraic closures, choose a
bijection of transcendence bases, and extend over the resulting
algebraic closures. Thus there are abstract field isomorphisms
\[
 \iota:\Omega_\ell\simeq\mathbb C,\qquad
 \iota':C_\ell\simeq\mathbb C,\qquad
           \iota j=\epsilon|_K,\quad\iota'j=\epsilon|_K.
                                                               \tag{NR.22}
\]
They need not agree off $j(K)$ and are not claimed continuous.
Only the explicitly displayed agreements are used.

\paragraph{NR13: the normalized characters and the sign of their reading.}
Retain \texttt{ValLL0} for the \emph{fixed} embedding
$\epsilon|_K$, the integrality of every exponent $e(j,g)$, and
the actual character frame $C$. HF13 constructs that frame;
HF16 supplies \texttt{ValLL0} only when its common value-field
generator is contained in the specified $K$. No enlargement of
an already fixed $D$ is implicit here.
Write $\delta_\ell=\delta'$ when $\ell=p_\delta$ and
$\delta_\ell=\delta$ otherwise. Its ramification set is the
primes over $p_\delta'$ or $p_\delta$, respectively, so it always
avoids $\ell$. For $1\le j\le m$ put
\[
 \beta_{\ell,j}
     =\mu_{0,j}\delta_\ell\eta^{-1}|\cdot|^{-r/2},
 \qquad
 \beta_{\ell,j}(P)
     =\mu_{0,j}(P)\delta_\ell(P)\eta(P)^{-1}N(P)^{r/2}.
                                                               \tag{NR.23}
\]
The actual $r$ is odd. The half-integral value law for $\eta$
and the integral value laws for $\mu_{0,j},\delta_\ell$ therefore
put these unramified values in $\epsilon(K)$, by GA11--GA12.

Take $K_0\supset c_0(\ell,j_0)$, its integer ring $O_0$, the
inclusion $\kappa:\operatorname{Frac}(O_0)=K_0\hookrightarrow
\Omega_\ell$, and $\iota$ from (NR.22). GA16 constructs the
actual continuous characters $\Delta:G_E\to O_0^\times$ and
$\chi_{1,j,b}:G_E\to O_0^\times$. The first is the finite
avatar of $\delta_\ell$; the second is the normalized avatar of
$\beta_{\ell,j}$, using $\iota$ for $b=\mathrm{true}$ and
$\overline{\phantom z}\circ\iota$ for $b=\mathrm{false}$.
Their coefficient embeddings are respectively $j_0$ and
$j_0\circ\bar{\phantom x}$. Conjugating twice shows that both
read the same specified complex embedding of $K$.

For the source's arithmetic Frobenius $\sigma_P$,
\[
 \iota\kappa(\Delta(\sigma_P))=\delta_\ell(P)^{-1},
 \qquad
 \iota_b\kappa(\chi_{1,j,b}(\sigma_P))
                                   =\beta_{\ell,j}(P)^{-1}
                                                               \tag{NR.24}
\]
at the permitted unramified primes. Since
\[
 (\mu_{0,j}\delta_\ell\eta^{-1})(P)
       =\beta_{\ell,j}(P)(\sqrt{N(P)})^{-r},            \tag{NR.25}
\]
the literal \texttt{NC.Read} exponent is $-r$, while that for
$\Delta$ is zero. To verify the record, discard the finite
ramification set and the primes over $\ell$, choose
$y\in K$ with $\epsilon(y)=\beta_{\ell,j}(P)^{-1}$, and apply
injectivity of $\iota_b\kappa$ to (NR.24). This gives the required
equality in $\operatorname{Frac}(O_0)$ for every chosen Frobenius
lift, with (NR.25) giving its complex value clause.
GA16 also proves the exact auxiliary inertia order $15$,
anticyclotomic law, off-set unramifiedness and prescribed tame
type-prime characters. They are properties of these normalized
avatars, not of freely chosen finite-order complex characters.
For the source's family indexed by every $j\in\mathbb N$, use
its actual reindexing $a(j)=\min(j-1,m-1)\in\{0,\ldots,m-1\}$,
where subtraction in this formula is truncated natural
subtraction and the rank data give $m\ge1$. Define the $j$th
character using row $a(j)$. This agrees with row $j-1$ on
$1\le j\le m$; all continuity, ramification and tame laws hold
also outside that range because they hold for every valid row.
The \texttt{Read} assertion itself is required only in the
literal range $1\le j\le m$.

\paragraph{NR14: construction of a genuinely nonempty twist-ring family.}
Fix once and for all $J_q$ from NR1. Given arbitrary families
$b_0,b_1$ of finite extensions with $b_0(\ell,j)\supset c_0(\ell,j)$,
and a finite set $S_x$ of primes of $E$, make the following choices
for every $q\in\operatorname{Spec}_1\mathcal O_K$:
\[
 \begin{split}
 \mathrm{Tw}&=\operatorname{Spec}_1\mathcal O_K,\qquad
 \ell(q)=\operatorname{char}\kappa(q),\qquad j_q=J_q,\\
 K_{0,q}&=b_0(\ell(q),j_q),\qquad
 L_q=b_1(\ell(q),j_q),\qquad O_{0,q}=\mathcal O_{K_{0,q}},\\
 O_q&=\mathcal O_{\widehat{K_{0,q}L_qU_{\ell(q)}}},\qquad
 O_{q,k}=\mathcal O_{\widehat{K_{0,q}L_qU_{\ell(q)}
                                      (\zeta_{\ell(q)^k})}}.
 \end{split}                                                   \tag{NR.26}
\]
The embedding $\mathrm{emb}_0:K\to\operatorname{Frac}(O_{0,q})$
is the codomain restriction of $J_q$, and
$\mathrm{emb}_{0b}(x)=\mathrm{emb}_0(\bar x)$. All algebra maps
between the displayed rings are their actual inclusions.
Use NR13 for $\Delta,\chi_1$ and NR11 for
$\operatorname{RCG}_{\ell(q),Q}$, its Artin quotient and its
restriction maps. Finally set
\[
 S_0(q)=\{v:v\mid\ell(q)\}\cup
       \bigcup_{a=0,1}\{v:v\mid g_a\}\cup
       \{v:v\mid p_\delta\}\cup\{v:v\mid p_\delta'\}\cup S_x.
                                                               \tag{NR.27}
\]
This set is finite because each nonzero rational integer has
only finitely many prime divisors in a number field.

These formulas supply every field of the source's \texttt{TwR}:
the finite $\mathbb Z_\ell$ coefficient DVR and finite residue
field are $O_{0,q}$; NR5--NR7 supply the complete DVR $O_q$,
algebraically closed residue field and finite cyclotomic tower;
NR13 supplies the actual normalized characters and exact
ramification clauses; NR11 supplies all the universal finite
groups and transition laws. To choose the conjugation of
$\Omega$, embed this algebraic closure of the number field into
$\mathbb C$. Its image is the field of algebraic numbers and
is stable under complex conjugation. Pulling conjugation back
gives an involution; on $E$ it is the host CM conjugation, since
every complex embedding of a CM field respects that conjugation.
This is the exact \texttt{IsCMConjOm} input used in NR10 and GA16.
Every rational prime has at least one prime above it in $K$:
the integral extension $\mathbb Z\subset\mathcal O_K$ has lying
over, or equivalently the finite ring $\mathcal O_K/\ell$ has a
maximal ideal. Hence this \texttt{Tw} is nonempty, indeed infinite.
For every one of its indices, $\Delta$ has a nontrivial inertia
image of order $15$.

\paragraph{NR15: all RingSpec compatibilities on the same fields.}
For the family (NR.26), use the identity equivalence
$\mathrm{Tw}\simeq\operatorname{Spec}_1\mathcal O_K$, the literal
identifications of $O_0,O,O_k$ with their integer rings, and the
bridges from NR3. Equation (NR.2) identifies the center of
$\mathrm{emb}_0$ with $q$ and proves \texttt{pl\_emb}.
The bounds in \texttt{RingSpec} hold with equality
\[
        K_{0,q}=b_0(\ell(q),j_q),\qquad
                   L_q=b_1(\ell(q),j_q).              \tag{NR.28}
\]
The diagrams for $\mathbb Z_\ell$, the base ring, and its finite
cyclotomic extensions commute because all maps are inclusions.
For example, for $z\ne0$, the two equations
$\mathrm{emb}_0(x)z=y$ and $J_q(x)z=y$ are equivalent after
the same injective fraction-field inclusion. The condition
\texttt{mu\_g} follows from the $\mu_\ell$ summand in (NR.21),
and the excluded set is exactly (NR.27).

Embed $\operatorname{Frac}(O_q)\subset C_{\ell(q)}$ and compose
with $\iota'$ from (NR.22) to define \texttt{iCq}. Define
$\mathrm{embO}(q,b)$ by composing $\mathrm{emb}_{0,b}$ with the
fraction-field inclusion. Then
\[
 \mathrm{iCq}_q\mathrm{embO}(q,\mathrm{true})=\epsilon|_K,\qquad
 \mathrm{embO}(q,\mathrm{false})(x)
                   =\mathrm{embO}(q,\mathrm{true})(\bar x).
                                                               \tag{NR.29}
\]
If a value of the base embedding is integral, this same inclusion
gives its required value in $O_q$, proving the sided fixed-field
clause without choosing another embedding.

Take $\mathrm{mu0H}(j)=\mu_{0,j}$ in its literal row range, and
the trivial character outside that range; take
$\mathrm{etaH}=\eta$, $\mathrm{delH}(q)=\delta_{\ell(q)}$.
NR13 proves the two \texttt{Read} clauses and their precise
exponents. HF11 gives the unramifiedness and the specified
unit-argument infinity law of $\eta$. The value equalities with
the original character frame hold by definition.
If two indices have the same auxiliary ramification set, their
$\mathrm{delH}$ agree. The only two candidate sets are the
nonempty sets of primes above the distinct rational primes
$p_\delta,p_\delta'$; these sets are disjoint, since a prime ideal
cannot contain two coprime rational primes. Thus equal sets
choose the same character. Finally the stronger
\texttt{deltaRead} assertion follows from the first equality
of (NR.24) at every permitted unramified prime, not merely
cofinitely many primes.

All fields of \texttt{RingSpec} have now been checked. We have
proved \texttt{exists\_ringSpec} with its exact order
\[
       \exists c_0\ \forall b_0\supset c_0\ \forall b_1\ \forall S_x
                   \ \exists(T,\mathrm{RingSpec}),     \tag{NR.30}
\]
and with the prescribed $J_q$ unchanged. Neither the rows nor
$K,L,\epsilon$ were replaced in this construction.

\paragraph{NR16: equality of centers determines the valuation ring.}
Let $v$ be a nontrivial additive valuation on a number field $K$,
with values in $\mathbb Z\cup\{\infty\}$. Every algebraic integer
has nonnegative value: it is integral over $\mathbb Z$, hence
over the valuation ring, which is integrally closed. Its center
\[
                   q_v=\{a\in\mathcal O_K:v(a)>0\}    \tag{NR.31}
\]
is a prime ideal. Additivity and the ultrametric inequality
give the ideal laws, and nonnegative values on integers show
that if $v(ab)>0$, one of $v(a),v(b)$ is positive.

Suppose $q_v=q\ne0$ and $e:K\hookrightarrow\operatorname{Frac}(O)$
is an embedding into a DVR whose maximal-ideal center is also $q$.
Algebraic integers map into $O$ by integral closedness. An element
$s\in\mathcal O_K\setminus q$ has $v(s)=0$ and $e(s)\in O^\times$.
Since $(\mathcal O_K)_q$ is a DVR, for every $x\in K$ either
$x=a/s$ or $x=s/a$, with $a,s\in\mathcal O_K$ and $s\notin q$.
In the first case $v(x)\ge0$ and $e(x)\in O$. In the second,
if $a\notin q$ both $x$ and $e(x)$ are units; if $a\in q$ then
$v(x)<0$ and $e(x)$ cannot be integral, because the product of
an integral element and $e(a)\in\mathfrak m_O$ cannot be the
unit $e(s)$. Therefore
\[
                 v(x)\ge0\quad\Longleftrightarrow\quad e(x)\in O
                               \qquad(x\in K).         \tag{NR.32}
\]
This upgrades a condition on integral elements of $K$ to the
condition on every element of $K$ required by \texttt{SidePl}.

\paragraph{NR17: exact all-place coefficient coverage.}
Let $w$ be any of the fixed places of $L$ in $D$. Restrict its
valuation to $K$ and take the center $q$ from (NR.31). It is
nonzero because the residue characteristic $p$ of $w$ has positive
value. NR15 gives the index $i=q$, whose coefficient embedding
has exactly this center. Both $p$ and $\ell(i)$ lie in $q$;
if they were different primes, a Bezout relation would put
$1$ in $q$. Hence $\ell(i)=p$. Applying NR16 yields
\[
 \forall x\in K,\qquad
   v_w(x)\ge0\quad\Longleftrightarrow\quad
                 \mathrm{emb}_0(i)(x)\in O_{0,i}.     \tag{NR.33}
\]
This is \texttt{exists\_setting\_of\_place}, and exactly the
$b=\mathrm{true}$ coefficient assertion in \texttt{SideB.cover}.
Taking positive values on $\mathcal O_K$ also proves the literal
\texttt{TwCover} center condition. The proof applies to every
$w$, irrespective of whether a row happens to be nonzero at $w$.
It does not assert the other fields of \texttt{SideB}.

\paragraph{NR18: a reading inducing a prescribed full-field place.}
The preceding coverage can be strengthened to the full
\texttt{Reading.Induces} condition, which concerns $L$, not just
$K$. Fix $w$ and the index $i$ from NR17. Choose
$k_w:L\hookrightarrow\Omega_{\ell(i)}$ inducing $w$ by NR1.
The restrictions $k_w|_K$ and $j_i$ induce the same valuation ring
on $K$. The rational-power comparison in NR2 proves that their
normalized absolute values agree. Their map on $K$ consequently
extends to an isomorphism of the finite local completions
\[
       \mathbb Q_\ell(k_w(K))\simeq\mathbb Q_\ell(j_i(K)),
                              \qquad k_w(x)\longmapsto j_i(x).
                                                               \tag{NR.34}
\]
An embedding of a finite separable extension into an algebraic
closure extends to an automorphism of that algebraic closure.
Choose such a $\mathbb Q_\ell$-automorphism $\tau$ and replace
$k_w$ by $\tau k_w$. Its norm is unchanged, so it still induces
$w$, and its restriction to $K$ is now exactly $j_i$.

Apply the transcendence-basis argument of NR12 to the number
field $L$ and the embeddings $k_w,\epsilon$. It gives an abstract
isomorphism $\iota_w:\Omega_\ell\simeq\mathbb C$ satisfying
$\iota_w k_w=\epsilon$ on all of $L$. Keep the canonical inclusion
$\kappa_i:\operatorname{Frac}(O_{0,i})\hookrightarrow\Omega_\ell$.
Then
\[
 \iota_w\kappa_i\mathrm{emb}_0(i)=\epsilon|_K,\qquad
 v_w(x)\ge0\ \Longleftrightarrow\
             |\iota_w^{-1}\epsilon(x)|\le1\quad(x\in L).
                                                               \tag{NR.35}
\]
The first equation and the unchanged $\mathbb Z_\ell$ inclusion
make $(\kappa_i,\iota_w)$ a literal \texttt{Reading}; the second
is its \texttt{Induces} clause. For a general already supplied
index and side, the identical argument applies whenever its
\texttt{SidePl} condition for $w$ holds, using
$\kappa_i\mathrm{emb}_{0,b}$ in place of $j_i$.
After any installation of this same ring family into an
$\mathrm{NTw}$, the assertion refers to that installation's
unchanged rings and embeddings. It supplies no signature,
residual representation or form comparison on its own.

\paragraph{NR19: a countable cofinal family of finite local fields.}
Fix an embedding $j:\overline{\mathbb Q}\hookrightarrow\Omega_\ell$.
For each $x\in\Omega_\ell$, choose $y\in j(\overline{\mathbb Q})$
so close to $x$ that
\[
 |x-y|<|x-x'|\quad\text{for every distinct conjugate }x'
                                  \text{ of }x\text{ over }\mathbb Q_\ell.
                                                               \tag{NR.36}
\]
There are finitely many such conjugates and NR2 gives density.
If $x\in\mathbb Q_\ell$, the conclusion below already holds.
Otherwise, every automorphism $\tau$ fixing $\mathbb Q_\ell(y)$
preserves the norm and satisfies
$|\tau x-x|\le\max(|\tau x-y|,|y-x|)=|x-y|$.
Equation (NR.36) forces $\tau x=x$. The fixed-field theorem in
the separable algebraic closure gives
\[
                              x\in\mathbb Q_\ell(y).  \tag{NR.37}
\]
This is the Krasner argument needed here, with its strict
inequality made explicit.

The field $\overline{\mathbb Q}$ is countable, and so is its set
of finite subsets. The fields
$K_S=\mathbb Q_\ell(j(S))$ for these finite subsets are finite
extensions. Every finite subset of $\Omega_\ell$ is contained
in one $K_S$, by applying (NR.37) to each of its elements.
Thus this is a countable cofinal family of finite subextensions,
although neither $\mathbb Q_\ell$ nor $\Omega_\ell$ is countable.

\paragraph{NR20: a compact representation has one finite coefficient field.}
Let $G$ be a compact Hausdorff topological group and
$\rho:G\to\operatorname{GL}_n(\Omega_\ell)$ a homomorphism whose
matrix entries are continuous. For the countable family of NR19
define
\[
 H_S=\{g:\text{all entries of }\rho(g)\text{ and }\rho(g^{-1})
                                                \text{ lie in }K_S\}.
                                                               \tag{NR.38}
\]
These are subgroups, by matrix multiplication and inversion.
They are closed: a finite extension of $\mathbb Q_\ell$ is
complete and therefore closed in $\Omega_\ell$, and there are
only finitely many continuous entry conditions. They cover
$G$, because a matrix and its inverse have finitely many entries.
The compact Hausdorff Baire theorem gives an $H_S$ with nonempty
interior. A subgroup with nonempty interior is open, and an open
subgroup of a compact group has finite index.

Choose representatives $g_1,\ldots,g_a$ for its finitely many
cosets. Adjoin their finitely many matrix entries to $K_S$,
obtaining a finite extension $K'$. For $g=g_jh$ with $h\in H_S$,
the equality $\rho(g)=\rho(g_j)\rho(h)$ puts every entry in $K'$.
The same assertion applied to $g^{-1}$ gives the inverse matrix
over $K'$. Hence the entire representation is defined over one
finite extension. If an initial finite field $K_{\mathrm{given}}$
must be retained, replace $K'$ by its compositum with that field.
This proves \texttt{exists\_finExt\_entries\_mem} with its arbitrary
lower-bound field. It does not assert that the image of $\rho$
is finite; compactness is used through Baire and boundedness,
not through such a claim.

\paragraph{NR21: a stable integral lattice and exact conjugation.}
With the finite field $K'$ from NR20, put $O'=\mathcal O_{K'}$
and choose a uniformizer $\pi$. The finitely many coordinate
functions of $\rho$ have bounded image on the compact group.
For some $N$,
$\rho(g)O'^n\subset\pi^{-N}O'^n$ for every $g$. Define
\[
                         \Lambda=\sum_{g\in G}\rho(g)O'^n,
              \qquad O'^n\subset\Lambda\subset\pi^{-N}O'^n.
                                                               \tag{NR.39}
\]
The sum denotes finite sums, so it is an $O'$-submodule.
The right-hand module is finite over the Noetherian DVR $O'$;
its submodule $\Lambda$ is therefore finite. It is torsion-free,
so the elementary divisor theorem over a DVR makes it free.
The two inclusions show that its rank is $n$. For every $h\in G$,
left multiplication permutes the summands, giving
$\rho(h)\Lambda=\Lambda$.

Choose an $O'$-basis of $\Lambda$ and let $P\in\operatorname{GL}_n(K')$
have those columns. Then
\[
                    P^{-1}\rho(g)P\in\operatorname{GL}_n(O')
                                  \quad\text{for all }g.       \tag{NR.40}
\]
The inverse is integral as well, since (NR.40) also holds for
$g^{-1}$. Matrix multiplication by the fixed $P,P^{-1}$ preserves
continuity. The induced topology on $O'$ is its adic topology,
so the resulting integral representation is continuous with
that topology. Conjugation preserves characteristic polynomials
and carries invariant direct-sum decompositions bijectively,
so it also preserves semisimplicity. This proves the actual
\texttt{exists\_finExt\_conj\_intOf} endpoint, including rank
zero if it occurs, without asserting that its integral model
has been identified with any pre-existing residual model.

\paragraph{NR22: the row twist has the exact Galois normalization.}
Now retain the literal \texttt{SideA} and \texttt{GalEx} of $D$.
The latter asserts, for every rational prime $\ell$, every complex
reading $\iota$, and every conjugate-self-dual regular algebraic
cuspidal $\pi$ of positive rank $a$, existence of a continuous
semisimple representation realizing $\pi$ at shift $(1-a)/2$.
This is a substantive automorphic Galois existence input, and
is not proved by NR20--NR21.

For $a=r$, let $\rho_1$ be the representation supplied by
\texttt{GalEx}. The actual odd integer $t$ and the half-integral
infinity type of $\eta$ make
\[
                     \gamma=\eta^{-1}|\cdot|^{-t/2}   \tag{NR.41}
\]
an algebraic Hecke character, as in GA11--GA13. Global Artin
reciprocity and GA6 give its continuous one-dimensional avatar
$\nu:G_E\to\Omega_\ell^\times$ with
$\iota(\nu(\sigma_P))=\gamma(P)^{-1}$ at the specified unramified
Frobenius elements. Define $\rho=\nu\rho_1$ by scalar multiplication.
It is continuous. A subspace is invariant for $\rho$ precisely
when it is invariant for $\rho_1$, so semisimplicity is preserved.

The \texttt{SideA} character reading identifies the local
one-dimensional representation $\eta_T^{-1}$ with $\eta^{-1}$
in precisely the source convention. The scalar
$u_P=\gamma(P)^{-1}$ is therefore the sole Frobenius eigenvalue
of that local character with the extra shift $-t/2$. If
$\alpha_1,\ldots,\alpha_r$ are the normalized eigenvalues of
$\rho_1(\sigma_P)$, then
\[
 \det(XI-u_P\,\iota(\rho_1(\sigma_P)))
                              =\prod_{a=1}^r(X-u_P\alpha_a),
 \qquad {1-r\over2}-{t\over2}={1-n_V\over2},
                          \quad n_V=r+t.              \tag{NR.42}
\]
Thus $\rho$ realizes $\pi\otimes\eta_T^{-1}$ at exactly
$(1-n_V)/2$. The finitely many ramified primes and primes over
$\ell$ can be removed together; the equality holds for every
Frobenius lift at each remaining prime. This proves
\texttt{exists\_isRepOf\_tw}, not a new assertion of \texttt{GalEx}.

\paragraph{NR23: keeping the representation while changing its reading.}
There are three precise invariances used in the bounded family.
First, suppose $\iota_0,\iota_1:\Omega_\ell\simeq\mathbb C$
both satisfy $\iota_bj=\epsilon|_K$. Suppose a representation
$\rho$ realizes $\pi$ under $\iota_0$, and that cofinitely many
of the required normalized Frobenius polynomials are
$p_P^\epsilon$ for polynomials $p_P\in K[X]$. Injectivity of
$\iota_0$ forces
\[
                \det(XI-\rho(\sigma_P))=p_P^j.
                                                               \tag{NR.43}
\]
Applying $\iota_1$ gives the same required complex polynomial.
Continuity and semisimplicity are properties of $\rho$ and
remain unchanged. No continuity of either $\iota_b$ is needed.

Second, if $\tau\in\operatorname{Aut}_{\mathbb Q_\ell}(\Omega_\ell)$,
it is an isometry, and applying it entrywise preserves continuity,
semisimplicity and matrix characteristic polynomials. Thus a
representation read with $\iota\tau$ becomes one read with
$\iota$ after applying $\tau$ to its coefficients.

Third, retain the literal \texttt{ChebSplit}: the specified
Frobenius lifts at split primes outside any finite set are dense
in $G_E$. If two continuous representations have the same
normalized Frobenius polynomials at those primes, every
coefficient of their characteristic polynomials agrees on a
dense subset. The coefficient functions are continuous into
the Hausdorff field $\Omega_\ell$, so
\[
                \det(XI-\rho(g))=\det(XI-\rho'(g))
                                      \quad(g\in G_E).           \tag{NR.44}
\]
In particular, if $\rho'$ realizes $\pi'$ and the local
representations of $\pi,\pi'$ agree at split primes, the original
$\rho$ also has the required Frobenius polynomials for $\pi'$.
This proves the source's \texttt{isRepOf\_of\_split}; it need not
choose an isomorphism between the two representations, and it
does not assert that split agreement constructs either one.

\paragraph{NR24: a uniform finite bound for all occurring row representations.}
In addition to the data of NR22, retain the actual
\texttt{RowCusp} and \texttt{ChebSplit}. Write
\[
 \mathrm{GoodOcc}(\sigma)\ \Longleftrightarrow\
 D.\mathrm{goodRow}(\sigma)\ \text{and}\
 \exists j\in[1,m],\ D.\mathrm{isot}(j,\sigma)\ne0.
                                                               \tag{NR.45}
\]
For each such $\sigma$, \texttt{RowCusp} supplies one actual
cuspidal $\pi_\sigma$, conjugate-self-dual and regular algebraic,
whose local representations agree with the given row at every
split prime. It also supplies, cofinitely in $P$, the
$K$-polynomials for $\pi_\sigma\otimes\eta_T^{-1}$ at the shift
$(1-n_V)/2$. This latter clause is what permits NR23; it is not
an automatic consequence of a name for the row.

For each $(\ell,j:K\hookrightarrow\Omega_\ell)$ and each good
occurring row, choose a complex reading extending $j$, apply
NR22, and apply NR21. Obtain an integral representation
$\rho_{\ell,j,\sigma}$ over a finite field $L_{\ell,j,\sigma}$.
NR23 shows that this \emph{same} representation works for every
complex reading extending $j$ and every conjugate-self-dual
regular algebraic $\pi$ with the same split local row. In the
latter comparison, \texttt{GalEx} first supplies a representation
for $\pi$; (NR.44) then transfers its Frobenius identities to the
already chosen representation.

There are finitely many row types $\sigma$ in the source. Put
$j_{\mathrm{true}}=j$ and $j_{\mathrm{false}}=j\circ\bar{\phantom x}$,
and define the finite extension
\[
 c_1(\ell,j)=\bigvee_{\sigma,\,b\in\{\mathrm{false},\mathrm{true}\}}
                               L_{\ell,j_b,\sigma}.   \tag{NR.46}
\]
Here $\bigvee$ denotes the compositum of this finite family;
for a row failing \texttt{GoodOcc}, use the base field and the
trivial representation, with no assertion of an automorphic
realization for that row. Equivalently, (NR.46) is the smallest
subfield of $\Omega_\ell$ containing all the displayed fields.
The bound $c_1$ is chosen before any $b_0,b_1,S_x,T$ or
$\mathrm{RingSpec}$.

For every supplied \texttt{RingSpec} with $b_0\supset c_1$,
its actual $O_{0,i}\simeq\mathcal O_{K_{0,i}}$ contains the
matrix entries of $\rho_{\ell(i),j_{i,b},\sigma}$. Restrict the
representation, including its inverse matrices, to this ring
and transport by that specified ring isomorphism. The result
is a continuous
\[
       \rho_{i,\sigma,b}:G_E\longrightarrow
                                      \operatorname{GL}_r(O_{0,i}).
                                                               \tag{NR.47}
\]
Continuity follows from the adic topology and its embedding
into $\Omega_{\ell(i)}$, as in GA1--GA3.

This family has the full universal reading clause of
\texttt{rho\_rows}. Namely, let
$\phi:\operatorname{Frac}(O_{0,i})\hookrightarrow\Omega_{\ell(i)}$
respect the specified $\mathbb Z_{\ell(i)}$ map, and let
$\iota\phi\mathrm{emb}_{0,b}=\epsilon|_K$.
The two embeddings $\phi$ and the canonical $\kappa_i$ of the
finite local field are conjugate over $\mathbb Q_{\ell(i)}$:
extend their isomorphism of finite subfields to obtain
$\phi=\tau\kappa_i$. Now
$\iota\tau j_{i,b}=\epsilon|_K$. The universal statement for the
chosen representation and the second part of NR23 show that
$\phi(\rho_{i,\sigma,b})$ realizes every permitted
$\pi\otimes\eta_T^{-1}$ at $(1-n_V)/2$. The equality is about
these matrices under these actual embeddings, not an unrelated
representation having the same dimension.

Finally choose $b_0(\ell,j)=c_0(\ell,j)c_1(\ell,j)$, or any
larger finite field, and any $b_1$. NR15 constructs the all-prime
ring family for these choices; the preceding construction then
gives its row representations. NR18 supplies readings inducing
each specified full-field place, and the universal reading
clause just proved applies to them as well. This is a coherent
simultaneous choice on the fixed $D$. Its only new coefficient
enlargements are the explicitly permitted local fields $b_0,b_1$;
the global fields in (NR.1) have not changed.

\medskip
\noindent\textbf{Exact remaining realization boundary.}
The complete source \texttt{SideB} contains more than coefficient
coverage and bounded Galois matrices. Its producer
\texttt{FormsStep.second\_step} also assumes the actual rational
row datum and integral Hecke lattice, $\mathrm{MtFam}$ and
$\mathrm{MtInput}$, the same level and tuple data, and the
automorphic inputs retained in NR22--NR24. A member
$\mathrm{MtReal}(\mu,j)$ is an injective map from the specified
space $D.\mathrm{Mt}(\mu,j)$ onto an entire adelic type space
$\mathrm{TypeSp}(\mathrm{wt},J,\chi)$, with an actual open level;
an arbitrary abstract vector-space isomorphism does not give
this assertion. The further $\mathrm{MtAd}$ ties that map
pointwise to the fixed $Y.\mathrm{IL}.jMt$, the same row frames,
the rational group and weight action, and the exact finite
matrix coordinates. None of these identities follows from
the coefficient constructions above.

Likewise $\mathrm{IsTripleWith}$ requires the actual occurring
signature and cuspidal parameter with its infinity and level
labels, a Galois representation for that parameter, the exact
scalar-twist equation for the already installed $\rho_S$, its
Hecke eigenvalue equations, and the specified residual
characteristic polynomials and Teichmuller reduction. NR18
supplies its \texttt{Induces} clause and NR22--NR24 supply the
stated conditional representation primitives, not these
remaining simultaneous identifications. $\mathrm{IsTripleFor}$
also requires $\mathrm{IsChars}$ for both readings, including
the already fixed tame-generator labels, and spectral coverage
of every relevant eigenspace. A compatible installation into
$\mathrm{NTw}$ still needs its actual form part $\mathrm{TwF}$
and toric part $\mathrm{TwL}$; nonempty local indices further
need $\mathrm{LocsCover}$ and the local form/Galois ties. The
bare stage with empty local indices does not furnish them.
Thus the nonempty all-prime family proved here closes concrete
arithmetic suppliers while leaving these exact same-object
geometric and automorphic interfaces explicit.
